\sectionkicker{Source-backed technical notes}
\subsection*{AI for Science — 専門化するモデル: Technical Notes}
本欄は記事本文で圧縮した一次資料上の情報を、比較・再検証しやすい形で整理したものである。時系列、確認済みの主張、留意点、一次資料URLを掲載する。
\medskip
\subsection*{Theme at a glance}
\addcontentsline{toc}{subsection}{Theme at a glance}
\begin{center}
\footnotesize
\begin{tabularx}{\linewidth}{@{}p{0.20\linewidth}p{0.10\linewidth}p{0.14\linewidth}X@{}}
\toprule
資料 & 位置づけ & 種別 & 時系列 \\
\midrule
GPT-Rosalind & 主要資料 & モデル & 2026-04-16 (モデル公開) \\
Decoding Scientific Experimental Images: The SPUR Benchmark for Perception, Understanding, and Reasoning & 補足資料 & 論文 & 2026-04-30 (論文公開) \\
Towards A Generative Protein Evolution Machine with DPLM-Evo & 補足資料 & 論文 & 2026-04-30 (論文公開) \\
\bottomrule
\end{tabularx}
\end{center}
\normalsize

\begin{technicalnote}{GPT-Rosalind}{主要資料}
\begingroup
% reader-facing Technical Notes break policy
\widowpenalty=10000
\clubpenalty=10000
\displaywidowpenalty=10000
\begin{tabularx}{\linewidth}{@{}>{\bfseries}p{0.22\linewidth}X@{}}
組織 & OpenAI \\
種別 & モデル \\
時系列 & 2026-04-16 (モデル公開) \\
\end{tabularx}
\smallskip
{\bfseries 一次資料から整理したtechnical points}
\begin{itemize}[leftmargin=1.5em,itemsep=0.35em]
\item \textbf{Vendor claim}: OpenAIはGPT-Rosalindを、drug discovery、genomics analysis、protein reasoning、scientific research workflowを対象とするfrontier reasoning modelとして紹介している。
\end{itemize}
\begin{minipage}{\linewidth}
% reader-facing Technical Notes generic boundary/source tail group
{\bfseries 読む際の境界}
\begin{itemize}[leftmargin=1.5em,itemsep=0.35em]
\item \textbf{分析上の留意点}: 公式一次資料から公開時期と記載された範囲は確認できるが、能力・安全性・性能に関する記述は独立再現された評価ではない。
\end{itemize}
\begin{samepage}
{\bfseries 一次資料}
\begingroup\sloppy
\begin{itemize}[leftmargin=1.5em,itemsep=0.25em]
\item {\scriptsize\url{https://openai.com/index/introducing-gpt-rosalind}}
\end{itemize}
\endgroup
\end{samepage}
\end{minipage}
\endgroup
\end{technicalnote}
\medskip

\begin{technicalnote}{Decoding Scientific Experimental Images: The SPUR Benchmark for Perception, Understanding, and Reasoning}{補足資料}
\begingroup
% reader-facing Technical Notes break policy
\widowpenalty=10000
\clubpenalty=10000
\displaywidowpenalty=10000
\begin{tabularx}{\linewidth}{@{}>{\bfseries}p{0.22\linewidth}X@{}}
組織 & - \\
種別 & 論文 \\
時系列 & 2026-04-30 (論文公開) \\
\end{tabularx}
\smallskip
{\bfseries 一次資料から整理したtechnical points}
\begin{minipage}{\linewidth}
% reader-facing Technical Notes coherent tail group
\begin{itemize}[leftmargin=1.5em,itemsep=0.35em]
\item \textbf{Author claim}: SPURは、科学実験画像のperception・understanding・reasoningを測るbenchmarkとして提案されている。
\end{itemize}
{\bfseries 読む際の境界}
\begin{itemize}[leftmargin=1.5em,itemsep=0.35em]
\item \textbf{分析上の留意点}: 結果や比較は著者報告で独立再現していない。保持している版は2026年5月26日更新だが、4月のchronologyは4月30日の初回投稿に基づくため、5月以降の本文変更を4月時点へ帰属させない。
\end{itemize}
\begin{samepage}
{\bfseries 一次資料}
\begingroup\sloppy
\begin{itemize}[leftmargin=1.5em,itemsep=0.25em]
\item {\scriptsize\url{http://arxiv.org/abs/2604.27604v2}}
\end{itemize}
\endgroup
\end{samepage}
\end{minipage}
\endgroup
\end{technicalnote}
\medskip

\begin{technicalnote}{Towards A Generative Protein Evolution Machine with DPLM-Evo}{補足資料}
\begingroup
% reader-facing Technical Notes break policy
\widowpenalty=10000
\clubpenalty=10000
\displaywidowpenalty=10000
\begin{tabularx}{\linewidth}{@{}>{\bfseries}p{0.22\linewidth}X@{}}
組織 & - \\
種別 & 論文 \\
時系列 & 2026-04-30 (論文公開) \\
\end{tabularx}
\smallskip
{\bfseries 一次資料から整理したtechnical points}
\begin{itemize}[leftmargin=1.5em,itemsep=0.35em]
\item \textbf{Author claim}: DPLM-Evoは、denoising中にsubstitution・insertion・deletionを明示的に予測するevolutionary discrete diffusion frameworkとして提案されている。
\end{itemize}
\begin{minipage}{\linewidth}
% reader-facing Technical Notes generic boundary/source tail group
{\bfseries 読む際の境界}
\begin{itemize}[leftmargin=1.5em,itemsep=0.35em]
\item \textbf{分析上の留意点}: 結果や比較は著者報告で独立再現していない。保持している版は2026年6月3日更新だが、4月のchronologyは4月30日の初回投稿に基づくため、5月以降の本文変更を4月時点へ帰属させない。
\end{itemize}
\begin{samepage}
{\bfseries 一次資料}
\begingroup\sloppy
\begin{itemize}[leftmargin=1.5em,itemsep=0.25em]
\item {\scriptsize\url{http://arxiv.org/abs/2605.00182v3}}
\end{itemize}
\endgroup
\end{samepage}
\end{minipage}
\endgroup
\end{technicalnote}
\medskip
